\documentclass[11pt,oneside]{book}

\usepackage{amsfonts,amsmath,amssymb,amsthm}
\usepackage{microtype}
\usepackage[hidelinks]{hyperref}
\hypersetup{
  pdftitle={The Mathematics Physics Demands},
  pdfauthor={Paul Gustafson}
}

\setlength{\emergencystretch}{3em}
\numberwithin{equation}{chapter}

\newtheorem{theorem}{Theorem}[chapter]
\newtheorem{proposition}[theorem]{Proposition}
\newtheorem{exercise}{Exercise}[chapter]

\newcommand{\C}{\mathbb C}
\newcommand{\R}{\mathbb R}
\newcommand{\Z}{\mathbb Z}
\newcommand{\End}{\operatorname{End}}
\newcommand{\Hom}{\operatorname{Hom}}
\newcommand{\im}{\operatorname{im}}
\newcommand{\tr}{\operatorname{tr}}
\newcommand{\id}{\operatorname{id}}
\newcommand{\physicalquestion}[1]{%
  \begin{quote}\textbf{Physical question.} #1\end{quote}}
\newcommand{\modelassumptions}[1]{%
  \begin{quote}\textbf{Model and assumptions.} #1\end{quote}}
\newcommand{\parttransition}[1]{%
  \begin{quote}\textbf{Why the next framework is necessary.} #1\end{quote}}

\title{The Mathematics Physics Demands\\
\large A Coordinate-Free Progression from Superposition to Anyons}
\author{Paul Gustafson}
\date{August 2026}

\begin{document}

\hypersetup{pageanchor=false}
\maketitle
\hypersetup{pageanchor=true}
\frontmatter
\pagestyle{plain}
\tableofcontents
\chapter*{Introduction}
\addcontentsline{toc}{chapter}{Introduction}

Why does physics repeatedly demand new mathematics?  A vibrating string first
asks only that possible motions can be added.  Quantum measurement then asks
for inner products and spectral projections.  Molecular spectra ask which
transformations preserve a system.  Electromagnetism asks how local data fit
together over spacetime.  Quantized Hall response asks which global features
survive deformation.  Open quantum systems ask that positivity remain valid
after arbitrary auxiliary systems are attached.  Anyons finally ask for a
language in which fusion and exchange are operations rather than afterthoughts.

This book follows those demands in the order
\[
\begin{gathered}
\text{superposition}\longrightarrow\text{symmetry}\longrightarrow
\text{fields}\\[-2pt]
\longrightarrow\text{topology}\longrightarrow\text{noncommutativity}
\longrightarrow\text{categorification}.
\end{gathered}
\]
The sequence is conceptual, not historical.  Physics decides why a structure
enters; mathematical prerequisites decide exactly where it can enter.

Each chapter begins with one physical question.  It then introduces the
smallest coordinate-free language that makes the question precise, proves one
central structural theorem, returns to the physical phenomenon, and ends with
exercises that connect the structure to models or experiments.  An
experimental observation is treated as data to be predicted or organized,
not as a proof of a mathematical formalism.  Every model states the
assumptions under which its conclusions are intended.

The companion volume \emph{Coordinate-free Mathematics} develops a broader
Algebra--Geometry--Analysis narrative.  The present book is not a rearranged
edition of it.  It has independent numbering, transitions, proofs, exercises,
and scope; the provenance ledger records the limited passages adapted from
the companion.


\mainmatter
\part{States, Superposition, and Measurement}

\chapter{Waves force linearity}
\label{ch:waves}

\physicalquestion{Why can two waves occupy the same string at the same time
without destroying one another?}

\modelassumptions{The string is homogeneous, perfectly flexible, held at
constant tension, and displaced transversely by a small amount.  Slopes are
small enough that quadratic geometric corrections, damping, and external
forcing may be neglected.}

\section{What language can express superposition?}

A real \emph{vector space} is a set $V$ with addition and scalar
multiplication satisfying the usual distributive, associative, and inverse
laws.  A map $T:V\to W$ is \emph{linear} when
\[
T(av+bw)=aT(v)+bT(w).
\]
These axioms are not introduced because coordinates are convenient.  They
encode the empirical modeling rule that amplitudes add and may be rescaled.

For a string of length $L$, let $u(x,t)$ be its transverse displacement.  In
the small-slope model, force balance gives
\begin{equation}
\rho\,\partial_t^2u=T\,\partial_x^2u,
\qquad u(0,t)=u(L,t)=0,
\label{eq:wave}
\end{equation}
where $T>0$ is tension and $\rho>0$ is linear density.

\section{What does the wave equation force?}

\begin{theorem}[Superposition]
The solutions of \eqref{eq:wave} form a real vector space.  For every fixed
time $t$, evolution from admissible initial data $(u_0,v_0)$ to
$(u(\cdot,t),\partial_tu(\cdot,t))$ is linear whenever the initial-value
problem has a unique solution.
\end{theorem}

\begin{proof}
Let $D=\rho\partial_t^2-T\partial_x^2$.  Differentiation is linear, so
$D(au+bv)=aD(u)+bD(v)$.  The boundary conditions are homogeneous and are
therefore preserved by the same combination.  Hence solutions are closed
under addition and scaling.  If $E_t$ denotes evolution, then
$aE_t(u_0,v_0)+bE_t(w_0,z_0)$ solves the equation with initial data
$(au_0+bw_0,av_0+bz_0)$.  Uniqueness identifies this solution with
$E_t(au_0+bw_0,av_0+bz_0)$.
\end{proof}

\section{What does this say about an observed string?}

The theorem predicts that independently prepared small-amplitude normal
modes pass through one another and that their measured displacement is the
sum of their displacements.  It does not predict exact superposition for a
large displacement: there the omitted nonlinear terms invalidate the model.
For fixed endpoints the functions
\[
u_n(x,t)=\sin\!\left(\frac{n\pi x}{L}\right)
\cos\!\left(\frac{n\pi}{L}\sqrt{T/\rho}\,t\right)
\]
are modes, and linearity makes every finite linear combination another
motion.  The coefficients are coordinates used to calculate a motion; the
solution space and its linear evolution exist before a basis is chosen.

\section{What further contact should be proved?}

\begin{exercise}[Energy distinguishes the physical solution space]
For sufficiently regular solutions of \eqref{eq:wave}, prove conservation of
\[
E(t)=\frac12\int_0^L\bigl(\rho|\partial_tu|^2
+T|\partial_xu|^2\bigr)\,dx.
\]
Explain which boundary term uses the fixed-end assumption.
\end{exercise}

\begin{exercise}[Critical damping breaks conservative mode evolution]
Write $m\ddot q+b\dot q+kq=0$ as a first-order linear system.  Classify its
three damping regimes by the characteristic polynomial and identify the
nontrivial Jordan block at $b^2=4mk$.
\end{exercise}


\chapter{Composite systems require tensor products}
\label{ch:tensors}

\physicalquestion{If two systems are prepared independently, what state
space contains both independent preparations and correlations between them?}

\modelassumptions{Only the linear structure of each component is retained.
The preparation rule is linear in either component when the other component
is fixed.}

\section{Why are pairs of vectors not enough?}

The direct sum $V\oplus W$ describes alternatives or independent sectors: a
vector is a pair $(v,w)$.  A quotient $V/N$ declares vectors differing by an
element of $N$ equivalent.  Neither construction represents a preparation
that is bilinear in two inputs.

The \emph{tensor product} is a vector space $V\otimes W$ with a bilinear map
$\tau:V\times W\to V\otimes W$ such that every bilinear
$b:V\times W\to X$ factors through a unique linear map
$\widetilde b:V\otimes W\to X$:
\[
b=\widetilde b\circ\tau.
\]
This universal property, rather than a choice of product basis, defines the
composite state space.

\section{What makes the tensor product canonical?}

\begin{theorem}[Existence and uniqueness of the tensor product]
Every pair of vector spaces has a tensor product.  Any two tensor products
satisfying the universal property are related by a unique isomorphism that
preserves pure tensors.
\end{theorem}

\begin{proof}
Take the free vector space on $V\times W$ and quotient by the span of the four
families enforcing additivity and homogeneity in each variable.  The class of
$(v,w)$ is denoted $v\otimes w$.  A bilinear map vanishes on those relations,
so its linear extension descends uniquely to the quotient.  If $(T,\tau)$ and
$(T',\tau')$ both satisfy the property, universality gives maps
$f:T\to T'$ and $g:T'\to T$ preserving pure tensors.  Uniqueness applied to
the identity factorization gives $gf=\id_T$ and $fg=\id_{T'}$.
\end{proof}

\section{Where does entanglement enter?}

A nonzero vector $\psi\in V\otimes W$ is a \emph{product vector} if
$\psi=v\otimes w$; otherwise it is entangled.  The universal property
permits all bilinear product preparations, but the resulting vector space
necessarily contains their linear combinations.  Thus entanglement is not
an added force or interaction.  It is already present in the minimal linear
space closed under superposition of composite preparations.

\section{What further contact should be proved?}

\begin{exercise}[Schmidt decomposition]
For finite-dimensional complex inner-product spaces $H_A,H_B$, associate to
$\psi\in H_A\otimes H_B$ the map $H_A^*\to H_B$.  Apply singular-value
decomposition to prove
$\psi=\sum_i s_i a_i\otimes b_i$ with orthonormal families and $s_i>0$.
Prove that $\psi$ is a product vector exactly when one Schmidt coefficient is
nonzero.
\end{exercise}

\begin{exercise}[Universal constructions]
Give coordinate-free universal properties for direct sums and quotients.
Use them to prove $k(X\sqcup Y)\cong kX\oplus kY$ and
$k(X\times Y)\cong kX\otimes kY$ for free vector spaces.
\end{exercise}


\chapter{Measurement requires Hilbert-space geometry}
\label{ch:measurement}

\physicalquestion{How can a state assign probabilities to mutually exclusive
measurement outcomes without choosing coordinates?}

\modelassumptions{The system is finite dimensional, pure states are rays in a
complex vector space, and an ideal repeatable measurement is represented by a
self-adjoint operator.}

\section{What structure turns amplitudes into probabilities?}

A Hermitian inner product on a complex vector space $H$ is positive definite,
linear in its second argument, and conjugate linear in its first.  It gives
$\|\psi\|^2=\langle\psi,\psi\rangle$.  The adjoint of $A:H\to H$ is the
unique $A^*$ satisfying
$\langle A\psi,\phi\rangle=\langle\psi,A^*\phi\rangle$.
An operator is self-adjoint when $A=A^*$ and normal when $A^*A=AA^*$.

An orthogonal projection is an operator $P=P^*=P^2$.  If $\|\psi\|=1$, the
number $\langle\psi,P\psi\rangle=\|P\psi\|^2$ lies in $[0,1]$ and is
independent of the phase chosen to represent the ray of $\psi$.

\section{Why do observables have spectral outcomes?}

\begin{theorem}[Finite-dimensional spectral theorem]
Every normal operator on a finite-dimensional complex inner-product space has
a unique decomposition
\[
A=\sum_{\lambda\in\sigma(A)}\lambda P_\lambda,
\qquad P_\lambda P_\mu=0\ (\lambda\ne\mu),
\qquad \sum_\lambda P_\lambda=1,
\]
where $P_\lambda$ projects onto the $\lambda$-eigenspace.
\end{theorem}

\begin{proof}
The characteristic polynomial gives an eigenvector $v$.  Normality implies
that $v^\perp$ is invariant under both $A$ and $A^*$: for $w\perp v$, use
$\ker(A-\lambda)^\perp=\im(A^*-\overline\lambda)$.  Induction on dimension
produces an orthonormal eigenbasis.  Grouping basis vectors by eigenvalue
gives the projections.  Their images are the intrinsic eigenspaces, so the
decomposition is unique.
\end{proof}

\section{How does the theorem return to measurement?}

For a self-adjoint observable $A$, all eigenvalues are real.  The ideal
measurement model assigns
\[
\Pr_\psi(A=\lambda)=\langle\psi,P_\lambda\psi\rangle,
\qquad
\sum_\lambda\Pr_\psi(A=\lambda)=1.
\]
The theorem supplies mutually exclusive outcomes and the inner product turns
their components into normalized probabilities.  This is a mathematical
model of measurement statistics; agreement of repeated experiments with
these frequencies tests the model rather than proving the spectral theorem.

\section{What further contact should be proved?}

\begin{exercise}[Uncertainty from noncommutativity]
For self-adjoint $A,B$ and unit $\psi$, set
$\Delta_\psi A=\|(A-\langle A\rangle_\psi)\psi\|$.  Use Cauchy--Schwarz to
prove
\[
\Delta_\psi A\,\Delta_\psi B
\geq\frac12|\langle\psi,[A,B]\psi\rangle|.
\]
State precisely when equality holds.
\end{exercise}

\begin{exercise}[Normal modes are spectral measurements]
For a finite chain of equal masses and springs, express its stiffness as a
positive self-adjoint operator.  Find its spectral projections and recover
the mode frequencies without treating the chosen mass coordinates as
intrinsic.
\end{exercise}


\chapter{Ancillas force complete positivity}
\label{ch:channels}

\physicalquestion{Why must a physical evolution remain positive after the
system is coupled to an arbitrary auxiliary system on which nothing is done?}

\modelassumptions{States of an $n$-level system are density operators in
$M_n(\C)$.  Dynamics is linear on statistical mixtures, and an unused
ancilla evolves by the identity map.}

\section{Why is positivity alone insufficient?}

A density operator is $\rho\geq0$ with $\tr\rho=1$.  A linear map
$\Phi:M_n\to M_m$ is positive when it carries positive matrices to positive
matrices.  It is \emph{completely positive} when
$\id_{M_r}\otimes\Phi$ is positive for every $r$.  A quantum channel is a
completely positive trace-preserving map.

The transpose is positive but $\id_{M_2}\otimes t$ sends the maximally
entangled two-qubit state to an operator with a negative eigenvalue.  Ordinary
positivity therefore fails the untouched-ancilla test.

For $\Phi:M_n\to M_m$, define the Choi matrix
\[
C_\Phi=\sum_{i,j}E_{ij}\otimes\Phi(E_{ij}).
\]

\section{What structure characterizes physical finite-dimensional maps?}

\begin{theorem}[Choi--Kraus--Stinespring]
For a linear map $\Phi:M_n\to M_m$, the following are equivalent:
\begin{enumerate}
\item $\Phi$ is completely positive;
\item $C_\Phi\geq0$;
\item $\Phi(x)=\sum_{a=1}^rK_axK_a^*$ for some $K_a:\C^n\to\C^m$.
\end{enumerate}
Consequently $\Phi(x)=V^*(x\otimes1_E)V$ for a map
$V:\C^m\to\C^n\otimes E$.  If $\Phi$ is unital, $V$ is an isometry.
\cite{choi,stinespring}
\end{theorem}

\begin{proof}
Complete positivity makes
$C_\Phi=(\id\otimes\Phi)(|\Omega\rangle\langle\Omega|)$ positive, with
$\Omega=\sum_i e_i\otimes e_i$.  A positive decomposition
$C_\Phi=\sum_a|v_a\rangle\langle v_a|$ identifies each $v_a$ with a matrix
$K_a$ and direct substitution gives the Kraus formula.  That formula is
positive after every matrix amplification.  Finally set
$V\xi=\sum_aK_a^*\xi\otimes e_a$; then
$V^*(x\otimes1)V=\sum_aK_axK_a^*$.  The identity
$V^*V=\Phi(1)$ proves the last assertion.
\end{proof}

\section{What does dilation say about open evolution?}

In the Schr\"odinger picture the adjoint form of the theorem says that every
channel can be written
\[
\rho\longmapsto\operatorname{Tr}_E(U(\rho\otimes\sigma)U^*)
\]
after enlarging the environment if necessary.  Irreversibility is therefore
compatible with reversible evolution of a larger closed system: information
is lost only when the environment is ignored.  Complete positivity is the
coordinate-free consistency condition that makes this statement survive
every added ancilla.

\section{What further contact should be proved?}

\begin{exercise}[No universal cloner]
Suppose a channel maps $\rho\otimes\sigma$ to $\rho\otimes\rho$ for every
pure state $\rho$.  Use monotonicity of fidelity, or contractivity of trace
distance, to derive a contradiction from two nonorthogonal states.
\end{exercise}

\begin{exercise}[Part I synthesis]
Begin with two normal modes of the string in Chapter~\ref{ch:waves}, regard
their complex span as a two-level system, attach a two-level environment, and
construct a unitary coupling whose reduced evolution is a dephasing channel.
Find its Kraus operators and Choi matrix, prove complete positivity, and
identify exactly which superpositions lose phase coherence.
\end{exercise}

\parttransition{Linear spaces, tensor products, and channels describe states
and evolutions.  They do not identify which transformations preserve the
physical structure or how those transformations label observable sectors.
That missing classification is symmetry.}

\part{Symmetry and Quantum Numbers}

\chapter{Symmetry separates state space into sectors}
\label{ch:representations}

\physicalquestion{Why do energy levels occur in symmetry-related families
rather than as unrelated individual states?}

\modelassumptions{A symmetry is an invertible transformation preserving the
physical predictions of the model.  The transformations compose as a finite
group $G$ and act linearly on a finite-dimensional complex state space $V$.}

\section{How does a physical transformation become algebra?}

A group representation is a homomorphism
$\rho:G\to\operatorname{GL}(V)$.  A subspace $W\subseteq V$ is invariant
when $\rho(g)W\subseteq W$ for all $g$.  It is irreducible when its only
invariant subspaces are $0$ and itself.  A map $T:V\to W$ intertwines two
representations when $T\rho_V(g)=\rho_W(g)T$.

If a Hamiltonian $H$ obeys $H\rho(g)=\rho(g)H$, then every eigenspace of $H$
is invariant.  Symmetry can therefore organize degeneracy only if invariant
spaces can be decomposed into elementary pieces.

\section{Why can finite symmetry be completely decomposed?}

\begin{theorem}[Maschke and Schur]
Every finite-dimensional complex representation of a finite group is a direct
sum of irreducible representations.  If $V,W$ are irreducible, every nonzero
intertwiner $V\to W$ is an isomorphism; moreover
$\End_G(V)=\C\,1_V$.
\end{theorem}

\begin{proof}
Average any Hermitian inner product over $G$ to make it invariant.  If
$W\subseteq V$ is invariant, then $W^\perp$ is invariant as well.  Induction
on dimension gives a decomposition into irreducibles.  For an intertwiner
$T$, its kernel and image are invariant, so irreducibility makes a nonzero
$T$ invertible.  If $T\in\End_G(V)$, choose an eigenvalue $\lambda$ over
$\C$.  The intertwiner $T-\lambda1$ has nonzero kernel and is therefore zero.
\end{proof}

\section{What does decomposition predict physically?}

Write
\[
V\cong\bigoplus_\alpha M_\alpha\otimes V_\alpha,
\]
where the $V_\alpha$ are inequivalent irreducibles and $G$ acts trivially on
the multiplicity spaces $M_\alpha$.  Every symmetry-preserving Hamiltonian
has the form $\bigoplus_\alpha H_\alpha\otimes1_{V_\alpha}$.  Thus an energy
inside the $\alpha$ sector has at least $\dim V_\alpha$-fold degeneracy unless
other structure splits the multiplicity space.  The claim is conditional on
the symmetry being exact; perturbations that break $G$ may split the level.

\section{What further contact should be proved?}

\begin{exercise}[Averaging produces invariant projectors]
For $W\subseteq V$ invariant and any projection $P:V\to W$, prove that
\[
\widetilde P=\frac1{|G|}\sum_{g\in G}\rho(g)P\rho(g)^{-1}
\]
is an equivariant projection onto $W$.
\end{exercise}

\begin{exercise}[Symmetry protects degeneracy]
Let an irreducible representation $V_\alpha$ occur once in an eigenspace of a
$G$-invariant Hamiltonian.  Prove that every $G$-invariant perturbation acts
there by a scalar.  Explain why this conclusion says nothing about a
perturbation that breaks $G$.
\end{exercise}


\chapter{Characters predict molecular spectra}
\label{ch:characters}

\physicalquestion{Why are some molecular vibrations visible in infrared or
Raman spectra while symmetry forbids others?}

\modelassumptions{The equilibrium molecule has a finite point group $G$.
Small nuclear displacements form a linear representation, and the electric
dipole approximation is valid for the transitions under discussion.}

\section{Which invariant detects an irreducible sector?}

The character of a representation is
$\chi_V(g)=\tr(\rho_V(g))$.  It is constant on conjugacy classes and does not
depend on a basis.  For class functions define
\[
\langle\chi,\psi\rangle_G
=\frac1{|G|}\sum_{g\in G}\overline{\chi(g)}\psi(g).
\]

For an irreducible $V_\alpha$ of dimension $d_\alpha$, the operator
\[
P_\alpha=\frac{d_\alpha}{|G|}
\sum_{g\in G}\overline{\chi_\alpha(g)}\rho_V(g)
\]
is the candidate projector onto its isotypic component.

\section{Why do characters recover multiplicities?}

\begin{theorem}[Character orthogonality]
Irreducible characters of a finite group are orthonormal.  If
$V\cong\bigoplus_\alpha m_\alpha V_\alpha$, then
\[
m_\alpha=\langle\chi_\alpha,\chi_V\rangle_G,
\]
and $P_\alpha$ is the projection onto the $V_\alpha$-isotypic component.
\end{theorem}

\begin{proof}
For a linear map $A:V_\alpha\to V_\beta$, average
$A\mapsto |G|^{-1}\sum_g\rho_\beta(g)A\rho_\alpha(g)^{-1}$.
The result is an intertwiner.  Schur's lemma from
Chapter~\ref{ch:representations} makes it zero for $\alpha\ne\beta$ and a
scalar identity for $\alpha=\beta$.  Applying this to matrix units and taking
traces gives character orthogonality.  Additivity of trace on direct sums
then gives the multiplicity formula.  Substitution into $P_\alpha$ shows that
it is identity on copies of $V_\alpha$ and zero on all other irreducibles.
\end{proof}

\section{How does this return to a water spectrum?}

For water at equilibrium the model point group is $C_{2v}$.  The nuclear
displacement representation has dimension nine.  Removing translations and
rotations leaves
\[
V_{\mathrm{vib}}\cong2A_1\oplus B_2.
\]
A dipole transition from $V_i$ to $V_f$ through a dipole component $V_\mu$
can be nonzero only if
\[
\Hom_G(1,V_f^*\otimes V_\mu\otimes V_i)\ne0.
\]
This invariant-tensor condition predicts allowed lines.  Measured line
positions additionally depend on masses, force constants, anharmonicity, and
instrumental conditions; symmetry organizes visibility but does not determine
the frequencies by itself.  Reference vibrational frequencies provide the
experimental comparison for this classification \cite{shimanouchi}.

\section{What further contact should be proved?}

\begin{exercise}[Water has three symmetry-classified vibrations]
Construct the $C_{2v}$ character of the nine-dimensional displacement space
from the fixed nuclei and spatial traces.  Subtract translations and rotations
and derive $2A_1\oplus B_2$.
\end{exercise}

\begin{exercise}[Raman selection differs from dipole selection]
Let the polarizability transform as $\operatorname{Sym}^2(V_\mu)$.  Determine
which $C_{2v}$ irreducibles occur and compare the Raman-active and
infrared-active water modes.
\end{exercise}


\chapter{Continuous symmetry produces angular momentum}
\label{ch:lie}

\physicalquestion{Why are rotational quantum numbers discrete even though
physical rotations vary continuously?}

\modelassumptions{Rotations act continuously and unitarily on a
finite-dimensional state space.  Infinitesimal generators are defined on the
whole space; domain issues of unbounded operators are postponed.}

\section{How is continuous transformation differentiated?}

A matrix Lie group is a closed subgroup of $\operatorname{GL}(n,\C)$.  Its
Lie algebra consists of the matrices $X$ for which $e^{tX}$ remains in the
group, with bracket $[X,Y]=XY-YX$.  Differentiating a smooth representation
$U:G\to U(H)$ gives a Lie-algebra representation.  For rotations, define
self-adjoint generators $J_i$ by differentiating the three one-parameter
rotation subgroups; with $\hbar=1$ they satisfy
\[
[J_i,J_j]=i\varepsilon_{ijk}J_k.
\]
Set $J^2=J_1^2+J_2^2+J_3^2$.

\section{Which operator labels rotational sectors?}

\begin{theorem}[The rotational Casimir]
The operator $J^2$ commutes with every $J_i$ and hence with the connected
rotation representation.  On every irreducible complex rotation sector it is
a scalar.
\end{theorem}

\begin{proof}
Using $[AB,C]=A[B,C]+[A,C]B$,
\[
[J^2,J_\ell]
=i\sum_{i,k}\varepsilon_{i\ell k}(J_iJ_k+J_kJ_i)=0,
\]
because the coefficient is antisymmetric in $i,k$ while the parenthesis is
symmetric.  Thus $J^2$ intertwines the Lie-algebra action and therefore the
connected group action.  Schur's lemma makes it scalar on an irreducible
sector.
\end{proof}

\section{Why do angular-momentum quantum numbers result?}

Introduce $J_\pm=J_1\pm iJ_2$.  The commutators
$[J_3,J_\pm]=\pm J_\pm$ make $J_\pm$ raising and lowering operators.  In a
finite-dimensional unitary irreducible sector, positivity of
$J_\mp J_\pm$ forces a highest weight $j\in\tfrac12\Z_{\ge0}$ and weights
$m=-j,-j+1,\ldots,j$, while $J^2=j(j+1)$.  Integer $j$ descends to an honest
representation of $SO(3)$; half-integer $j$ is an honest representation of
its double cover $SU(2)$ and only a projective representation of $SO(3)$.
The continuous group therefore produces discrete labels through unitarity and
finite dimensionality.

\section{What further contact should be proved?}

\begin{exercise}[Spin one-half]
Set $J_i=\sigma_i/2$ for the Pauli matrices.  Verify the commutation
relations, compute $J^2$, and prove that a $2\pi$ rotation acts by $-1$ while
a $4\pi$ rotation acts by $1$.
\end{exercise}

\begin{exercise}[Rotational invariance conserves angular momentum]
Let canonical position and momentum operators satisfy
$[Q_i,P_j]=i\delta_{ij}$ and set
$L_i=\sum_{jk}\varepsilon_{ijk}Q_jP_k$.  For
$H=(P_1^2+P_2^2+P_3^2)/(2m)+V(Q_1^2+Q_2^2+Q_3^2)$, prove
$[H,L_i]=0$ and identify precisely where rotational invariance of the
potential is used.
\end{exercise}


\chapter{Quantum phases require projective symmetry}
\label{ch:projective}

\physicalquestion{How can a symmetry act consistently on physical rays even
when its operators fail to multiply exactly as the group does?}

\modelassumptions{State vectors differing by a phase represent the same pure
state.  A discrete symmetry group acts unitarily on rays, and phases belong to
$U(1)$ with trivial group action.}

\section{Which phases are removable?}

A projective representation has operators satisfying
\[
U_gU_h=\alpha(g,h)U_{gh},\qquad \alpha(g,h)\in U(1).
\]
Associativity forces the cocycle equation
\[
\alpha(g,h)\alpha(gh,k)=\alpha(h,k)\alpha(g,hk).
\]
Rephasing $U_g$ by $\lambda(g)$ changes $\alpha$ by the coboundary
$\lambda(g)\lambda(h)\lambda(gh)^{-1}$.  The quotient of cocycles by
coboundaries is $H^2(G,U(1))$.

\section{Why does cohomology classify projective actions?}

\begin{theorem}[Projective classification]
Equivalence classes of normalized factor systems for projective unitary
representations of $G$ are classified by $H^2(G,U(1))$.  Tensor product
multiplies the corresponding cohomology classes.
\end{theorem}

\begin{proof}
The displayed associativity calculation says exactly that a factor system is
a normalized $2$-cocycle.  Changing representatives of rays rephases the
operators and hence multiplies the cocycle by a coboundary; conversely every
coboundary is produced by such a rephasing.  For tensor products,
$(U_g\otimes V_g)(U_h\otimes V_h)=
\alpha(g,h)\beta(g,h)U_{gh}\otimes V_{gh}$, so classes multiply.
\end{proof}

\section{How can a cohomology class protect an edge state?}

Let $G=\Z/2\times\Z/2$ be generated by two $\pi$ rotations.  On a spin-half
edge take $U_x=\sigma_1$ and $U_z=\sigma_3$.  Their anticommutation realizes
the nontrivial element of $H^2(G,U(1))\cong\Z/2$.  Any local edge operator
commuting with both is scalar, so a symmetry-preserving local perturbation
cannot split the doublet.  In the spin-one Haldane model, the statement relies
on a bulk gap, locality, and preservation of the specified symmetry; closing
the gap or breaking the symmetry evades the conclusion.

\section{What further contact should be proved?}

\begin{exercise}[The projective class squares to zero]
Compute the factor system for $U_x,U_z$ above and prove that its tensor square
is cohomologically trivial.  Interpret this as the possibility of combining
two protected edges into an ordinary linear representation.
\end{exercise}

\begin{exercise}[Part II synthesis]
Let a Hamiltonian commute with a finite symmetry group and suppose one
isotypic sector carries a projective edge action.  Use characters to isolate
the sector, Schur's lemma to constrain symmetric perturbations, and
$H^2(G,U(1))$ to decide whether the action can be rephased to an ordinary
representation.  Separate degeneracy forced by an irreducible dimension from
degeneracy protected by a nontrivial projective class.
\end{exercise}

\parttransition{Global symmetry uses the same group at every point.  A field
must compare locally varying values and frames across spacetime.  Making
symmetry local forces differential forms, metrics, bundles, connections, and
curvature.}

\part{Classical Fields and Geometry}

\chapter{Maxwell fields are differential forms}
\label{ch:forms}

\physicalquestion{How can electric charge and magnetic flux be described in
a way that every observer can integrate consistently?}

\modelassumptions{Spacetime is a smooth oriented four-manifold.  The
electromagnetic field is classical, sources are smooth on the region under
study, and singular charges or material interfaces are excluded unless they
are represented by removed subsets or distributions.}

\section{What can be pulled back and integrated?}

A smooth $n$-manifold is a Hausdorff, second-countable space locally modeled
on $\R^n$ by smoothly compatible charts.  Its tangent and cotangent bundles
are $TM$ and $T^*M$.  A differential $p$-form is a smooth section of
$\Lambda^pT^*M$; write $\Omega^p(M)$ for their space.

The wedge product is alternating and the exterior derivative
$d:\Omega^p(M)\to\Omega^{p+1}(M)$ is the unique graded derivation extending
the differential of functions and satisfying $d^2=0$.  For a smooth map
$f:N\to M$, pullback obeys $f^*d=df^*$.  Integration pairs compactly
supported top-degree forms with oriented manifolds.  Stokes' theorem states
\[
\int_Xd\omega=\int_{\partial X}\omega.
\]

Closed forms lie in $\ker d$, exact forms in $\im d$, and
$H^p_{\mathrm{dR}}(M)=\ker d/\im d$ records the obstruction to finding a
global potential.

\section{Which conservation law is structural?}

\begin{theorem}[Maxwell charge conservation]
Let $F\in\Omega^2(M)$ be the electromagnetic field and let
$J\in\Omega^3(M)$ be the current.  If a metric-dependent Hodge star satisfies
the Maxwell equations
\[
dF=0,\qquad d{\star F}=J,
\]
then $dJ=0$.  For every compact oriented four-region $X$,
\[
\int_{\partial X}J=0.
\]
\end{theorem}

\begin{proof}
Apply $d$ to the sourced equation: $dJ=d^2{\star F}=0$.  Stokes' theorem then
gives $\int_{\partial X}J=\int_XdJ=0$.
\end{proof}

\section{How does this return to charge measurements?}

The boundary of a spacetime slab consists of two spatial slices and a timelike
side.  The theorem says that the difference of total charge on the slices
equals the current flux through the side.  This conclusion does not depend on
a coordinate continuity equation; that equation is a chart expression of
$dJ=0$.  The physical model still assumes that $J$ includes every relevant
source.  Apparent nonconservation may instead signal current crossing an
unmodeled boundary or failure of the classical description.

\section{What further contact should be proved?}

\begin{exercise}[Flux depends only on homology]
If $F$ is closed and two oriented $2$-cycles differ by the boundary of a
$3$-chain, prove that their integrals of $F$ agree.  State the dual dependence
on the de Rham class $[F]$.
\end{exercise}

\begin{exercise}[Gauge potentials are local data]
If $F=dA$ on a contractible region, prove that replacing $A$ by $A+d\lambda$
does not change $F$.  Give an example on a punctured region where a closed
one-form is not exact.
\end{exercise}


\chapter{Metrics turn geometry into dynamics}
\label{ch:metric}

\physicalquestion{What additional structure turns invariant field forms into
lengths, wave equations, freely falling paths, and tidal acceleration?}

\modelassumptions{The manifold carries a smooth nondegenerate metric of fixed
signature.  Test bodies are small enough to neglect their backreaction and,
when geodesic motion is invoked, non-gravitational forces are absent.}

\section{What does a metric determine without coordinates?}

A pseudo-Riemannian metric is a smooth nondegenerate symmetric bilinear form
$g$ on each tangent space.  An orientation and $g$ determine a volume form
$\operatorname{vol}_g$ and a Hodge star by
\[
\alpha\wedge\star\beta
=\langle\alpha,\beta\rangle_g\operatorname{vol}_g.
\]

A connection on $TM$ assigns covariant derivatives $\nabla_XY$, is linear in
$X$, satisfies a Leibniz rule in $Y$, and permits vectors at nearby points to
be compared infinitesimally.  It is torsion free when
$\nabla_XY-\nabla_YX=[X,Y]$ and metric compatible when $\nabla g=0$.

\section{Why is there a canonical metric connection?}

\begin{theorem}[Levi--Civita]
Every pseudo-Riemannian metric has a unique torsion-free, metric-compatible
connection.  It is determined by the coordinate-free Koszul formula
\begin{align*}
2g(\nabla_XY,Z)={}&Xg(Y,Z)+Yg(Z,X)-Zg(X,Y)\\
&-g(X,[Y,Z])+g(Y,[Z,X])+g(Z,[X,Y]).
\end{align*}
\end{theorem}

\begin{proof}
Metric compatibility expands each derivative of an inner product into two
covariant-derivative terms.  Add the expansions for $Xg(Y,Z)$ and
$Yg(Z,X)$, subtract the one for $Zg(X,Y)$, and use zero torsion to replace
differences of covariant derivatives by Lie brackets.  This yields the
formula and hence uniqueness by nondegeneracy of $g$.  Conversely the formula
defines $\nabla_XY$ uniquely; direct substitution verifies linearity, the
Leibniz rule, zero torsion, and metric compatibility.
\end{proof}

\section{How do connection and star return to physics?}

The Hodge star in Chapter~\ref{ch:forms} supplies the metric part of Maxwell's
equations.  The Levi--Civita connection defines geodesics by
$\nabla_{\dot\gamma}\dot\gamma=0$ and curvature by
\[
R(X,Y)Z=\nabla_X\nabla_YZ-\nabla_Y\nabla_XZ-\nabla_{[X,Y]}Z.
\]
For a variation of geodesics with separation field $J$, the equation
$\nabla_T\nabla_TJ=R(T,J)T$ turns curvature into relative acceleration.  A
gravitational-wave interferometer uses this tidal effect; it does not measure
coordinate components of the metric in isolation.

\section{What further contact should be proved?}

\begin{exercise}[The Hodge star is unique]
Prove pointwise existence and uniqueness of $\star$ from the induced inner
product on exterior powers.  Determine $\star^2$ in terms of form degree and
metric signature.
\end{exercise}

\begin{exercise}[Killing fields give conserved quantities]
If $X$ satisfies $\mathcal L_Xg=0$ and $\gamma$ is an affinely parametrized
geodesic, prove that $g(X,\dot\gamma)$ is constant.
\end{exercise}


\chapter{Hamiltonian evolution is symplectic}
\label{ch:symplectic}

\physicalquestion{Why does conservative time evolution preserve phase-space
volume even when it stretches and folds individual regions?}

\modelassumptions{The classical system has a finite-dimensional smooth phase
space, its dynamics is deterministic, and the Hamiltonian has no explicit
time dependence.  Dissipation and stochastic forcing are excluded.}

\section{What replaces a metric on phase space?}

A symplectic form is a closed nondegenerate two-form $\omega$ on a manifold
$M$.  Nondegeneracy identifies every one-form with a unique vector field.  A
Hamiltonian $H\in C^\infty(M)$ therefore determines $X_H$ by
\[
\iota_{X_H}\omega=dH.
\]
The Poisson bracket is $\{f,h\}=\omega(X_f,X_h)$.  In canonical coordinates
the definition reproduces Hamilton's equations, but neither $\omega$ nor
$X_H$ depends on those coordinates.

\section{Why can Hamiltonian flow not compress phase volume?}

\begin{theorem}[Liouville]
The local flow $\varphi_t$ of a Hamiltonian vector field preserves the
symplectic form and the Liouville volume $\omega^n/n!$ on a $2n$-dimensional
phase space.
\end{theorem}

\begin{proof}
Cartan's formula and $d\omega=0$ give
\[
\mathcal L_{X_H}\omega
=d(\iota_{X_H}\omega)+\iota_{X_H}d\omega=d^2H=0.
\]
Thus
$\frac d{dt}\varphi_t^*\omega
=\varphi_t^*\mathcal L_{X_H}\omega=0$,
so $\varphi_t^*\omega=\omega$.  Pullback respects wedge products, hence it
also fixes $\omega^n/n!$.
\end{proof}

\section{What physical consequence follows?}

An ensemble transported by a Hamiltonian flow cannot converge into a smaller
phase-volume attractor.  This explains why ordinary friction cannot be
represented by an autonomous Hamiltonian on the same finite-dimensional
phase space: friction contracts volumes and violates the theorem's
assumptions.  In equilibrium statistical mechanics the preserved volume is
the invariant measure from which microcanonical counting begins.

\section{What further contact should be proved?}

\begin{exercise}[Magnetic curvature gives the Lorentz force]
On $T^*Q$ replace the canonical symplectic form by
$\omega_B=\omega_0+q\pi^*B$ for a closed two-form $B$ on $Q$.  For kinetic
Hamiltonian $H$, derive the Lorentz-force equation.
\end{exercise}

\begin{exercise}[Moment maps recover angular momentum]
For the lifted rotation action on $T^*\R^3$, prove that the moment map is
$J(q,p)=q\times p$.  Show directly that a rotation-invariant Hamiltonian
conserves $J$, and compare with Chapter~\ref{ch:lie}.
\end{exercise}


\chapter{Local symmetry becomes a connection}
\label{ch:gauge}

\physicalquestion{How can internal states at different spacetime points be
compared when no preferred frame exists globally?}

\modelassumptions{Internal states form fibers of a smooth vector bundle.  A
compact Lie group $G$ acts on the fibers, local gauge transformations are
changes of frame, and classical gauge fields are smooth connections.}

\section{Which object compares neighboring fibers?}

A vector bundle $E\to M$ is locally a product $U\times V$ with linear
transition maps.  A connection is a linear map
\[
\nabla:\Omega^0(M,E)\to\Omega^1(M,E),
\qquad \nabla(fs)=df\otimes s+f\nabla s.
\]
Extending it to $E$-valued forms defines curvature
$F_\nabla=\nabla^2\in\Omega^2(M,\End E)$.

For a principal $G$-bundle, a local frame represents the connection by a
Lie-algebra-valued one-form $A$, with
\[
F_A=dA+A\wedge A.
\]
A gauge transformation changes $A$ but conjugates $F_A$; curvature is the
local obstruction to path-independent transport.

\section{Which identity constrains every gauge field?}

\begin{theorem}[Bianchi identity]
Every connection satisfies $d_A F_A=0$, or intrinsically
$\nabla F_\nabla=0$.
\end{theorem}

\begin{proof}
In a local frame, $d_A\eta=d\eta+A\wedge\eta-(-1)^p\eta\wedge A$ for a
$p$-form.  Substitute $F_A=dA+A\wedge A$, use $d^2=0$ and the graded Leibniz
rule, and observe that all remaining terms cancel associatively.  Since the
identity transforms covariantly, it is independent of the chosen frame.
\end{proof}

\section{How does local symmetry produce field equations?}

With an invariant inner product on the Lie algebra and the Hodge star from
Chapter~\ref{ch:metric}, the Yang--Mills action is
\[
S(A)=\frac12\int_M\langle F_A\wedge\star F_A\rangle.
\]
A compactly supported variation gives $d_A\star F_A=0$ in vacuum.  The Bianchi
identity is kinematic; the Yang--Mills equation is dynamical.  For $G=U(1)$
they reduce to the two Maxwell equations.  Nonabelian groups additionally
self-couple through $A\wedge A$.

The Standard Model uses a particular gauge group and matter representations.
Those choices are empirical inputs to the model; gauge geometry organizes
their charges and couplings but does not derive the observed representation
list from differential geometry alone.

\section{What further contact should be proved?}

\begin{exercise}[Gauge variation of curvature]
For $A^g=g^{-1}Ag+g^{-1}dg$, prove
$F_{A^g}=g^{-1}F_Ag$ and show that the Yang--Mills action is gauge invariant.
\end{exercise}

\begin{exercise}[Part III synthesis]
On a principal $U(1)$-bundle, connect the four descriptions
\[
\text{potential}\longrightarrow\text{connection}\longrightarrow
\text{curvature}\longrightarrow\text{Maxwell field}.
\]
Derive both Maxwell equations, distinguish the Bianchi identity from the
Euler--Lagrange equation, and identify exactly where orientation, metric, and
bundle topology enter.
\end{exercise}

\parttransition{Curvature and differential equations are local.  They cannot
detect every effect of a flat connection around a hole, nor explain why a
transport coefficient remains an integer under continuous deformation.
Global observables force holonomy, characteristic classes, K-theory, and
index.}

\part{Topological Quantum Phenomena}

\chapter{A field-free loop can carry phase}
\label{ch:ab}

\physicalquestion{How can an electron acquire a measurable phase while
traveling only through a region where the electromagnetic field vanishes?}

\modelassumptions{An ideal solenoid confines magnetic flux to an inaccessible
region.  Outside it the particle is coherent, nonrelativistic, and described
by minimal $U(1)$ coupling.  Leakage fields and environmental decoherence are
neglected in the ideal model.}

\section{What does the interference observation demand?}

Two electron paths passing on opposite sides of a shielded solenoid recombine.
The model predicts a relative phase even though $F=dA=0$ along each path.
The gauge-invariant datum is not the value of $A$ at a point but the holonomy
\[
\operatorname{Hol}_A(\gamma)
=\exp\!\left(\frac{iq}{\hbar}\int_\gamma A\right)
\]
around the closed loop obtained from the two paths.

In the field-free region $M$, the potential is a closed one-form.  A gauge
change $A\mapsto A+d\lambda$ changes its representative but not its de Rham
class $[A]\in H^1_{\mathrm{dR}}(M)$.

\section{Why is holonomy topological in the flat region?}

\begin{theorem}[Cohomology--homology dependence]
For a closed one-form $A$ and a closed loop $\gamma$, the integral
$\int_\gamma A$ depends only on $[A]\in H^1_{\mathrm{dR}}(M)$ and
$[\gamma]\in H_1(M;\R)$.  Consequently the Aharonov--Bohm phase is invariant
under gauge change and homology-preserving deformation of the path.
\end{theorem}

\begin{proof}
Replacing $A$ by $A+d\lambda$ changes the integral by
$\int_\gamma d\lambda=0$.  If $\gamma-\gamma'=\partial C$, then Stokes'
theorem from Chapter~\ref{ch:forms} gives
\[
\int_\gamma A-\int_{\gamma'}A=\int_CdA=0.
\]
Both dependencies therefore descend to the stated classes.
\end{proof}

\section{What phase does the ideal solenoid predict?}

For $M=\R^3$ minus the solenoid axis, let
\[
A=\frac{\Phi}{2\pi}d\theta.
\]
A once-winding loop has $\int_\gamma A=\Phi$, so the relative phase is
$q\Phi/\hbar$ modulo $2\pi$.  Shielded-flux electron interferometry measures
the displacement of fringes as flux varies.  The experiment tests this
holonomy prediction while separately controlling leakage fields; it does not
prove de Rham theory \cite{aharonov-bohm,tonomura}.

\section{What further contact should be proved?}

\begin{exercise}[Winding computes holonomy]
For a loop of winding number $n$ around the removed axis, prove
$\int_\gamma d\theta=2\pi n$ and compute its phase.
\end{exercise}

\begin{exercise}[Flux periodicity]
Show that all interference probabilities are unchanged when
$q\Phi/\hbar$ changes by $2\pi k$.  Identify the corresponding flux quantum.
\end{exercise}


\chapter{Curvature produces stable characteristic classes}
\label{ch:chern}

\physicalquestion{Which global quantity detects a twisted family of quantum
states even though every small patch admits a smooth frame?}

\modelassumptions{The parameter space is a compact smooth manifold and the
relevant eigenspaces have constant finite rank.  Connections are smooth and
the structure group is unitary.}

\section{How can local curvature define a global class?}

For a complex vector bundle $E\to X$ with connection $\nabla$, let
$F_\nabla$ be its curvature.  The total Chern form is defined by
\[
\det\!\left(1+\frac{i}{2\pi}F_\nabla\right)
=1+c_1(E,\nabla)+c_2(E,\nabla)+\cdots.
\]
The forms are assembled from invariant polynomials, so they agree under
changes of local frame.

Stable equivalence declares $E$ and $F$ equivalent when
$E\oplus\underline\C^r\cong F\oplus\underline\C^s$ after balancing ranks.
The Grothendieck group of bundles under direct sum is $K^0(X)$.

\section{Why do Chern classes ignore the chosen connection?}

\begin{theorem}[Chern--Weil independence]
Each Chern form is closed, and its de Rham class is independent of the
connection.  Direct sum makes the Chern character
\[
\operatorname{ch}(E)=\left[\tr\exp\!\left(\frac{iF_\nabla}{2\pi}\right)\right]
\]
descend to a homomorphism
$\operatorname{ch}:K^0(X)\to H^{\mathrm{even}}_{\mathrm{dR}}(X)$.
\end{theorem}

\begin{proof}
The Bianchi identity from Chapter~\ref{ch:gauge} and invariance of trace imply
$d\,\tr(F_\nabla^k)=0$.  For a path of connections $\nabla_t$ with derivative
$a_t$, differentiating under the trace gives
\[
\frac d{dt}\tr(F_t^k)=k\,d\,\tr(a_t\wedge F_t^{k-1}).
\]
Integrating in $t$ shows that two choices differ by an exact form.  The
exponential trace is additive on direct sums, so it respects the Grothendieck
relation and descends to $K^0(X)$.
\end{proof}

\section{What invariant results for a surface?}

For a complex bundle over a closed oriented surface,
\[
C_1(E)=\frac{i}{2\pi}\int_X\tr(F_\nabla)
\]
is an integer.  It is unchanged by a smooth deformation of the connection or
by adding trivial bundles.  Curvature may shift from place to place, but its
total Chern number cannot vary continuously.  The integer can jump only when
the underlying bundle ceases to vary continuously within the assumed class.

\section{What further contact should be proved?}

\begin{exercise}[A nonzero Chern number obstructs a global frame]
Prove that a bundle admitting a global frame is trivial and has vanishing
positive-degree Chern classes.  Conclude that $C_1(E)\ne0$ obstructs a global
smooth eigenbasis.
\end{exercise}

\begin{exercise}[Stable equivalence preserves the Chern character]
Show directly that adding a trivial bundle changes only the degree-zero rank
term of the Chern character.
\end{exercise}


\chapter{The Hall integer is a bundle invariant}
\label{ch:hall}

\physicalquestion{Why does Hall conductance remain pinned to integer
multiples of $e^2/h$ while microscopic Hamiltonian parameters vary?}

\modelassumptions{The crystal is noninteracting and translation invariant,
the Fermi energy lies in a spectral gap, temperature is negligible compared
with that gap, and the occupied bands vary smoothly over the Brillouin torus.
The TKNN linear-response formula is taken as the physical bridge from the
model to conductance.}

\section{How does a gapped Hamiltonian produce topology?}

Let $H:k\mapsto H(k)$ be a smooth family of finite-dimensional self-adjoint
operators over the Brillouin torus $\mathbb T^2$.  A uniform gap at zero gives
the occupied spectral projection
\[
P(k)=\frac{1}{2\pi i}\int_\Gamma(z-H(k))^{-1}\,dz,
\]
where $\Gamma$ surrounds the negative spectrum.  The images of $P(k)$ form
the Bloch bundle $E_{\mathrm{occ}}\to\mathbb T^2$.  Adding inert filled or
empty flat bands changes the bundle by a trivial summand, so the stable phase
is a K-theory class.

\section{Why can the Hall integer change only at a gap closing?}

\begin{theorem}[Gap-preserving invariance]
A smooth gap-preserving homotopy $H_s(k)$ gives isomorphic endpoint occupied
bundles.  Hence
\[
\mathrm{Ch}_1(P)
=\frac{i}{2\pi}\int_{\mathbb T^2}
\tr\bigl(P\,dP\wedge dP\bigr)
\]
is an integer independent of $s$ and of added trivial bands.
\end{theorem}

\begin{proof}
Uniform compactness of $[0,1]\times\mathbb T^2$ keeps a common contour inside
the gap, so the Riesz formula produces a smooth projection on the product.
Its image is a vector bundle whose restrictions at $s=0$ and $s=1$ are
isomorphic because an interval is contractible.  The displayed form
represents the first Chern class by Chapter~\ref{ch:chern}; its integral is
therefore constant and integral.  A trivial summand has zero curvature and
does not change the pairing.
\end{proof}

\section{How does topology return to conductance?}

The TKNN formula predicts
\[
\sigma_{xy}=\mathrm{Ch}_1(P)\frac{e^2}{h}.
\]
This is the physical bridge from the bundle invariant to the observed
response \cite{tknn,von-klitzing}.
Thus continuous changes of hopping amplitudes cannot move the ideal
conductance away from its plateau while the gap and model assumptions remain
valid.  A jump requires a point where the occupied projection is no longer
defined continuously---typically a gap closing.  Disorder and interactions
require broader formulations, but the invariant-versus-gap logic survives in
operator-algebraic versions.

\section{What further contact should be proved?}

\begin{exercise}[Berry curvature equals projection curvature]
For an isolated occupied line with local unit vector $\psi(k)$, compare
$i\langle\psi,d\psi\rangle$ with the connection induced by $P$.  Prove that
their curvatures yield the projection formula above.
\end{exercise}

\begin{exercise}[A two-band transition]
For $H(k)=d(k)\cdot\sigma$ with $d(k)\ne0$, show that the occupied Chern
number is the degree of $d/|d|:\mathbb T^2\to S^2$.  Identify why the degree
can change only where $d=0$.
\end{exercise}


\chapter{Fredholm index counts protected zero modes}
\label{ch:index}

\physicalquestion{Why is the imbalance of two kinds of zero modes stable
even though each kernel dimension can jump under perturbation?}

\modelassumptions{Operators act between complex Hilbert spaces, have closed
range, and have finite-dimensional kernel and cokernel.  Perturbations are
bounded and small in norm, or compact when explicitly stated.}

\section{Which difference survives perturbation?}

An operator $T:H_0\to H_1$ is Fredholm when its kernel and cokernel are finite
dimensional and its range is closed.  Its index is
\[
\operatorname{ind}T=\dim\ker T-\dim\operatorname{coker}T.
\]
The compact operators $\mathcal K(H)$ form an ideal in the bounded operators,
and the Calkin algebra is $\mathcal Q(H)=\mathcal B(H)/\mathcal K(H)$.

\section{Why is the index topologically stable?}

\begin{theorem}[Atkinson stability]
$T$ is Fredholm exactly when its image in the Calkin algebra is invertible.
The Fredholm index is locally constant and is unchanged by compact
perturbations.
\end{theorem}

\begin{proof}
If $T$ is Fredholm, choose inverses on complements of its finite-dimensional
kernel and cokernel; extending by zero gives a parametrix $S$ for which
$ST-1$ and $TS-1$ have finite rank.  Conversely a parametrix modulo compacts
implies finite-dimensional kernel and cokernel and closed range.  Thus
Fredholm operators are the inverse image of the invertible group of
$\mathcal Q(H)$, an open set.  Along a sufficiently small norm path the
kernel-cokernel difference cannot change, as follows from the parametrix and
finite-dimensional reduction.  Compact perturbations have the same Calkin
image and can be joined through $T+sK$, so the index is unchanged.
\end{proof}

\section{How does index meet geometry and physics?}

An elliptic differential operator on a closed manifold has a Fredholm
extension.  The Atiyah--Singer theorem identifies its analytic index with a
topological K-theory pairing of its principal symbol \cite{atiyah-singer}.
In gauge backgrounds,
this converts characteristic charge into chiral zero-mode imbalance.  The
deep theorem is used here as a bridge; the stability proved above already
explains why admissible perturbations cannot change the imbalance.

For the circle operator $D=d/d\theta$ on periodic complex functions, kernel
and cokernel are both the constants, so $\operatorname{ind}D=0$, agreeing
with the Euler characteristic of $S^1$.  This elementary case displays the
analytic-topological equality without claiming a proof of the general index
theorem.

\section{What further contact should be proved?}

\begin{exercise}[The unilateral shift has index minus one]
On $\ell^2(\mathbb N)$ let $S(e_n)=e_{n+1}$.  Compute the kernels and
cokernels of $S$ and $S^*$, then verify index stability after a finite-rank
perturbation.
\end{exercise}

\begin{exercise}[Part IV synthesis]
For the Aharonov--Bohm phase, Hall integer, and Fredholm index, identify the
structured input, the deformation class, the numerical pairing, and the
hypothesis whose failure permits a jump.  Express the first as a
cohomology--homology pairing, the second as a Chern pairing of a K-class, and
the third as the boundary invariant of an invertible Calkin class.
\end{exercise}

\parttransition{Topology explains quantized sectors but not how maps behave
after arbitrary matrix amplification, how nonlocal correlations are bounded,
or how infinite observable algebras contain physical weak limits.  Those
questions require noncommutative analysis.}

\part{Quantum Information and Noncommutative Analysis}

\chapter{Composite tests require matrix norms}
\label{ch:operator-spaces}

\physicalquestion{Why can two linear maps with the same ordinary norm behave
differently when an entangled auxiliary system is attached?}

\modelassumptions{Systems are finite dimensional in the explicit
calculations.  A subspace of operators inherits the norm arising from its
action on a Hilbert space, and an $r$-level ancilla is represented by matrix
amplification.}

\section{Which norm remembers every ancilla?}

A concrete operator space is a closed subspace
$E\subseteq\mathcal B(H,K)$.  Its $r$th matrix level is
\[
M_r(E)\subseteq\mathcal B(H^{\oplus r},K^{\oplus r}).
\]
For $\phi:E\to F$, set $\phi_r([x_{ij}])=[\phi(x_{ij})]$ and
\[
\|\phi\|_{\mathrm{cb}}=\sup_r\|\phi_r\|.
\]
The row and column operator spaces are
$R_n=M_{1,n}$ and $C_n=M_{n,1}$ with their concrete matrix norms.  As Hilbert
spaces both are $\C^n$; as operator spaces they encode different directions
of action.

\section{Can ordinary norm miss an arbitrarily large effect?}

\begin{theorem}[Transposition detects matrix levels]
For transposition $t:M_n\to M_n$,
\[
\|t\|=1,
\qquad
\|t\|_{\mathrm{cb}}=n.
\]
Thus ordinary boundedness does not control behavior under ancillas.
\end{theorem}

\begin{proof}
Transposition preserves singular values, so its operator norm on $M_n$ is
one.  At matrix level $n$, apply $\id\otimes t$ to the flip
$F=\sum_{ij}E_{ij}\otimes E_{ji}$, which is unitary and has norm one.  The
image is
$\sum_{ij}E_{ij}\otimes E_{ij}=n|\Omega\rangle\langle\Omega|$ for normalized
$\Omega=n^{-1/2}\sum_i e_i\otimes e_i$, and has norm $n$.  Hence the cb norm
is at least $n$.  Writing a matrix amplification as block transposition and
using the row--column factorization gives the reverse bound $n$.
\end{proof}

\section{What does this mean for composite quantum systems?}

The transpose preserves positivity of an isolated density matrix but, as
seen in Chapter~\ref{ch:channels}, fails positivity on an entangled ancilla.
The norm calculation measures the same obstruction quantitatively: the gap
between bounded and completely bounded behavior grows with dimension.  Matrix
levels are therefore operational data, not decorative copies of a Banach
norm.

\section{What further contact should be proved?}

\begin{exercise}[Rows and columns are completely different]
Prove that the identity maps between the underlying Hilbert spaces of $R_n$
and $C_n$ have cb norm $\sqrt n$ in both directions.
\end{exercise}

\begin{exercise}[Ruan's concrete identities]
For $x\in M_m(E)$ and $y\in M_n(E)$ prove
\[
\|x\oplus y\|=\max\{\|x\|,\|y\|\},\qquad
\|\alpha x\beta\|\leq\|\alpha\|\|x\|\|\beta\|.
\]
Explain why these identities are invariant under complete isometry.
\end{exercise}


\chapter{Noncommutative products require Haagerup tensoring}
\label{ch:haagerup}

\physicalquestion{Which tensor product records the norm of a process formed
by multiplying or composing noncommuting operator-valued inputs?}

\modelassumptions{Operator spaces are concrete and finite dimensional for
examples, although the definitions extend by completion.  Bilinear maps are
tested at all matrix levels using matrix multiplication of indices.}

\section{Why does the algebraic tensor product forget too much?}

For $u=\sum_{k=1}^r x_k\otimes y_k\in E\otimes F$, define
\[
\|u\|_h=\inf
\left\{
\|[x_1\ \cdots\ x_r]\|_{M_{1,r}(E)}
\|[y_1\ \cdots\ y_r]^{\mathsf T}\|_{M_{r,1}(F)}
\right\}.
\]
The completion is the Haagerup tensor product $E\otimes_hF$.

For a bilinear map $b:E\times F\to G$, define
\[
b_n(x,y)_{ij}=\sum_kb(x_{ik},y_{kj}),
\qquad \|b\|_{\mathrm{mb}}=\sup_n\|b_n\|.
\]

\section{Which universal property fits multiplication?}

\begin{theorem}[Haagerup linearization]
A bilinear map $b:E\times F\to G$ is multiplicatively bounded exactly when
its linearization extends to a completely bounded map
\[
\widetilde b:E\otimes_hF\to G.
\]
Moreover $\|\widetilde b\|_{\mathrm{cb}}=\|b\|_{\mathrm{mb}}$.
\end{theorem}

\begin{proof}
A factorization $u=a\odot c$ with
$a\in M_{1,r}(E)$ and $c\in M_{r,1}(F)$ is sent at matrix level to the
corresponding product $b_n(a,c)$.  Hence
$\|\widetilde b_n(u)\|\leq\|b\|_{\mathrm{mb}}\|a\|\|c\|$; taking the
infimum proves one inequality and extension by continuity.  Conversely apply
$\widetilde b_n$ to elementary matrix factorizations $x\odot y$ to obtain
$\|b_n(x,y)\|\leq\|\widetilde b\|_{\mathrm{cb}}\|x\|\|y\|$.  Taking
suprema gives equality.
\end{proof}

\section{Where does this tensor norm appear physically?}

Multiplication in a $C^*$-algebra is multiplicatively bounded and therefore
linearizes through $\otimes_h$.  The complete isometries
\[
C_n\otimes_hR_n\cong M_n,
\qquad
R_n\otimes_hC_n\cong S_1^n
\]
distinguish operator norm from trace norm solely by reversing factors.  The
second norm controls optimal state discrimination, so the order-sensitive
tensor geometry has an operational meaning.

\section{What further contact should be proved?}

\begin{exercise}[Helstrom discrimination]
For states $\rho_0,\rho_1$ with equal prior probability, prove that the best
binary measurement succeeds with probability
\[
\frac12+\frac14\|\rho_0-\rho_1\|_1.
\]
Relate the trace norm to $R_n\otimes_hC_n$.
\end{exercise}

\begin{exercise}[Forgetting matrix norms loses tensor information]
Show that $R_n$ and $C_n$ have isometric underlying Banach spaces while their
Haagerup products in the two orders have different norms.  Explain why no
symmetric completely isometric flip can identify the two products.
\end{exercise}


\chapter{Channels are distinguished at the matrix level}
\label{ch:diamond}

\physicalquestion{How well can one use distinguish two quantum processes when
the input may be entangled with an auxiliary system?}

\modelassumptions{Channels act between finite matrix algebras, are used once,
and have equal prior probability.  Arbitrary finite-dimensional ancillas and
joint output measurements are allowed.}

\section{Which distance includes every composite probe?}

Recall complete positivity and channels from Chapter~\ref{ch:channels}.  For
a linear map $\Psi:M_n\to M_m$, its diamond norm is
\[
\|\Psi\|_\diamond
=\sup_r\|\Psi\otimes\id_{M_r}\|_{1\to1}.
\]
The supremum explicitly asks how distinguishability changes when an untouched
ancilla is attached.  Operator systems---self-adjoint unital subspaces of
$\mathcal B(H)$---supply the matrix order that makes complete positivity
intrinsic even when the domain is not a full algebra.

\section{Why is the diamond norm operational?}

\begin{theorem}[One-use channel discrimination]
For equally likely channels $\Phi_0,\Phi_1:M_n\to M_m$, the optimal one-use
success probability is
\[
p_{\mathrm{succ}}
=\frac12+\frac14\|\Phi_0-\Phi_1\|_\diamond.
\]
The ancillary dimension $n$ is sufficient.
\end{theorem}

\begin{proof}
For an input state $\rho$ on system plus ancilla, the two output states differ
by $((\Phi_0-\Phi_1)\otimes\id)(\rho)$.  The Helstrom formula from
Chapter~\ref{ch:haagerup} gives the optimal measurement for that input.
Maximizing over inputs produces the diamond norm.  Convexity reduces the
maximum to pure inputs, and Schmidt decomposition from Chapter~\ref{ch:tensors}
shows that a pure input needs Schmidt rank at most $n$; a larger ancilla cannot
improve the value.
\end{proof}

\section{What advantage can entanglement provide?}

If the maximizing input is entangled, a strategy restricted to uncorrelated
probes attains a smaller norm.  The advantage is not additional channel use;
it comes from evaluating the same map at a higher matrix level.  Process
discrimination experiments compare observed success frequencies with the
separable benchmark and the diamond-norm prediction under calibrated noise
assumptions.

\section{What further contact should be proved?}

\begin{exercise}[Two unitary channels]
For unitary channels $\rho\mapsto U_i\rho U_i^*$, express perfect one-use
distinguishability in terms of the numerical range of $U_0^*U_1$.
\end{exercise}

\begin{exercise}[The Heisenberg adjoint has cb norm]
Prove in finite dimensions that
$\|\Psi\|_\diamond=\|\Psi^*\|_{\mathrm{cb}}$ when the adjoint is equipped
with operator norm matrix levels.
\end{exercise}


\chapter{Bell correlations are noncommutative geometry}
\label{ch:bell}

\physicalquestion{How strongly can separated quantum measurements correlate
without allowing communication between the laboratories?}

\modelassumptions{Each party chooses one of two binary observables.  The
observables of different parties commute in the tensor-product model, each
outcome is $\pm1$, and the experimental settings are statistically independent
of the prepared state.}

\section{Which operator encodes the Bell test?}

For self-adjoint unitaries $A_0,A_1$ and $B_0,B_1$, define the CHSH operator
\[
S=A_0\otimes(B_0+B_1)+A_1\otimes(B_0-B_1).
\]
Local hidden-variable assignments satisfy $|\langle S\rangle|\leq2$.  The
operators need not commute within one laboratory, and that noncommutativity
changes the bound.

\section{How large can the quantum value be?}

\begin{theorem}[Tsirelson bound]
Every quantum realization above satisfies
$\|S\|\leq2\sqrt2$, and the bound is attained by a maximally entangled pair
of qubits with suitable Pauli measurements \cite{tsirelson}.
\end{theorem}

\begin{proof}
Using $A_i^2=B_j^2=1$, expand to obtain
\[
S^2=4\,1-[A_0,A_1]\otimes[B_0,B_1].
\]
Each commutator has norm at most two, so $\|S^2\|\leq8$ and
$\|S\|\leq2\sqrt2$.  For the singlet state, take orthogonal Pauli directions
for Alice and the normalized sum and difference directions for Bob; direct
calculation gives magnitude $2\sqrt2$.
\end{proof}

\section{What does a violation establish experimentally?}

Under the stated locality and statistical assumptions, an observed value
above two rules out the corresponding local hidden-variable model.  It does
not permit signaling, since each local marginal is independent of the remote
choice.  Grothendieck-type inequalities generalize the theorem: for broad
families of correlation experiments, quantum advantage over classical vector
assignments remains bounded by a universal constant.  Loophole-controlled
Bell experiments compare measured correlations with these bounds under
explicitly tested locality assumptions \cite{hensen}.

\section{What further contact should be proved?}

\begin{exercise}[Every pure entangled two-qubit state violates CHSH]
Put a pure two-qubit state in Schmidt form and optimize planar Pauli
measurements to obtain a CHSH value strictly above two whenever both Schmidt
coefficients are nonzero.
\end{exercise}

\begin{exercise}[Grothendieck turns operators into vectors]
For a real coefficient matrix $(c_{ij})$, formulate the classical sign
optimization and the quantum vector relaxation.  Show how Tsirelson's vector
representation places their ratio under the real Grothendieck constant.
\end{exercise}


\chapter{Infinite observables lead to subfactor index}
\label{ch:subfactors}

\physicalquestion{How can one measure the relative size of nested infinite
observable algebras when ordinary vector-space dimension is meaningless?}

\modelassumptions{Observable algebras act nondegenerately on Hilbert space and
are closed in the weak operator topology.  For the index calculation, the
algebras are finite factors with faithful normalized trace.}

\section{Which closure contains physical limits?}

A von Neumann algebra is a unital $*$-subalgebra
$M\subseteq\mathcal B(H)$ equal to its bicommutant $M''$.  Equivalently, it
is weak-operator closed.  Its predual selects the normal functionals, the
states compatible with monotone limits of observables.

For an inclusion $N\subseteq M$ of finite factors, a trace-preserving
conditional expectation $E_N:M\to N$ is the noncommutative analogue of
averaging onto the smaller observable algebra.  On $L^2(M)$ let $e_N$ be the
orthogonal projection onto $L^2(N)$.

\section{Why do nested expectations produce Temperley--Lieb relations?}

\begin{theorem}[Jones projections]
In the iterated basic construction for a finite-index inclusion, the Jones
projections $e_i$ satisfy
\[
e_i^2=e_i=e_i^*,\qquad e_ie_j=e_je_i\ (|i-j|\ge2),
\qquad e_ie_{i\pm1}e_i=\delta^{-2}e_i,
\]
where $\delta^2=[M:N]$ is the Jones index.
\end{theorem}

\begin{proof}
Idempotence and self-adjointness hold because each $e_i$ is an orthogonal
projection.  Projections separated by at least two stages act on commuting
levels of the tower.  For adjacent stages, use
$e_Nxe_N=E_N(x)e_N$ and the Markov normalization of the trace in the basic
construction; applying the expectation twice contributes the scalar
$[M:N]^{-1}=\delta^{-2}$.  Density of the algebraic tower extends the
relations to $L^2$.
\end{proof}

\section{How does index return to quantum structure?}

The parameter $\delta$ behaves like the dimension of a bimodule and composes
multiplicatively.  Positivity of the Temperley--Lieb trace forces the values
below two to be $\delta=2\cos(\pi/n)$, producing the discrete Jones indices
$4\cos^2(\pi/n)$ \cite{jones-index}.  The same relations yield braid-group representations;
what began as relative size of observable algebras becomes data for particle
exchange.  The quantization theorem is deeper than the displayed relations
and is used here as the bridge to Part VI.

\section{What further contact should be proved?}

\begin{exercise}[Conditional expectation gives a Jones projection]
Prove $e_Nxe_N=E_N(x)e_N$ first on the dense subspace $M\subseteq L^2(M)$.
\end{exercise}

\begin{exercise}[Part V synthesis]
Track one composite-system question through four levels: ordinary norm,
matrix norm, completely positive order, and von Neumann closure.  For each
level identify a physical operation it controls and a counterexample showing
why the preceding level is insufficient.  End by deriving the braid
generators from the Temperley--Lieb projections.
\end{exercise}

\parttransition{Operator spaces and algebras encode states, maps, and
inclusions, but they do not make fusion, reassociation, exchange, defects, and
gluing into first-class operations.  Those processes demand monoidal and
braided categories.}

\part{Anyons, Topology, and Generalized Theories}

\chapter{Planar exchange produces braid groups}
\label{ch:braids}

\physicalquestion{Why can exchanging identical particles in two dimensions
retain information that disappears in three dimensions?}

\modelassumptions{Particles are pointlike, identical, and constrained to a
two-dimensional region.  Collision configurations are excluded, motion is
adiabatic, and worldlines are compared up to deformation that fixes their
endpoints.}

\section{Which group records planar worldlines?}

For $n$ ordered particles in a surface $\Sigma$, remove collisions from
$\Sigma^n$ and quotient by permutations to obtain the configuration space of
identical particles.  Its fundamental group is the braid group $B_n$ when
$\Sigma$ is a disk.  It has generators $\sigma_1,\ldots,\sigma_{n-1}$ and
relations
\[
\sigma_i\sigma_{i+1}\sigma_i
=\sigma_{i+1}\sigma_i\sigma_{i+1},
\qquad
\sigma_i\sigma_j=\sigma_j\sigma_i\quad(|i-j|\ge2).
\]
Unlike the symmetric group, $B_n$ does not impose $\sigma_i^2=1$.

A monoidal category has a tensor product, a unit, and coherent associativity
maps.  A braiding is a natural family
$c_{x,y}:x\otimes y\to y\otimes x$ satisfying the two hexagon identities.

\section{Why does categorical braiding act on many-particle states?}

\begin{theorem}[Braiding represents the braid group]
For any object $x$ in a braided monoidal category, the maps
\[
\rho(\sigma_i)
=1_x^{\otimes(i-1)}\otimes c_{x,x}\otimes
1_x^{\otimes(n-i-1)}
\]
define a representation of $B_n$ on $x^{\otimes n}$, with associators
inserted coherently.
\end{theorem}

\begin{proof}
Naturality and the hexagon identities give the Yang--Baxter equality on three
adjacent tensor factors, which is the first braid relation.  Maps acting on
disjoint tensor factors commute, giving the second.  Mac Lane coherence makes
the result independent of how the iterated tensor product is parenthesized.
\end{proof}

\section{What physical information survives an exchange?}

Adiabatic transport around a braid may act on the degenerate state space by
$\rho(\beta)$.  In an abelian anyon sector this action is a phase; in a
nonabelian sector it can be a noncommuting matrix.  Three-dimensional exchange
adds a homotopy that collapses the braid to a permutation, forcing the more
familiar boson/fermion alternatives under standard assumptions.  Planar
topology removes that collapse and therefore demands braid data.

\section{What further contact should be proved?}

\begin{exercise}[The symmetric group is a quotient]
Prove that imposing $\sigma_i^2=1$ on $B_n$ gives the symmetric group $S_n$.
Interpret the added relation as forgetting overcrossing versus undercrossing.
\end{exercise}

\begin{exercise}[Temperley--Lieb generators braid]
Given the relations in Chapter~\ref{ch:subfactors}, determine scalars $a,b$
for which $a1+be_i$ satisfy the braid relations.
\end{exercise}


\chapter{Fusion requires tensor categories}
\label{ch:fusion}

\physicalquestion{When anyons are combined, why can the result have several
possible particle types and an internal fusion space?}

\modelassumptions{There are finitely many topological charge types, fusion
spaces are finite dimensional, every charge has an antiparticle, and local
microscopic excitations have been quotiented out.}

\section{Which language contains fusion and exchange together?}

In a semisimple rigid $\C$-linear monoidal category, simple objects represent
anyon types.  Fusion is the decomposition
\[
x\otimes y\cong\bigoplus_z N_{xy}^z z,
\qquad
\Hom(z,x\otimes y)
\]
is the corresponding fusion space.  Rigidity supplies an antiparticle $x^*$.
An associator compares the two orders of fusing three particles, while a
braiding exchanges them.  Pentagon and hexagon coherence express consistency
of all composite processes.

The monodromy $M_{x,y}=c_{y,x}c_{x,y}$ is a full winding of $x$ around $y$.

\section{How is abelian monodromy classified?}

\begin{theorem}[Pointed monodromy]
Suppose every simple object is invertible and the simple classes form the
cyclic group $\Z/m$ generated by $a$.  A scalar monodromy compatible with
fusion has the form
\[
M_{a^r,a^s}=\exp\!\left(\frac{2\pi i krs}{m}\right)
\]
for some $k\in\Z/m$.
\end{theorem}

\begin{proof}
The hexagon identities make monodromy multiplicative in either variable, so
$M$ is a bicharacter on $\Z/m\times\Z/m$.  It is determined by
$\zeta=M_{a,a}$.  Since $a^m\cong1$, multiplicativity gives $\zeta^m=1$;
write $\zeta=e^{2\pi ik/m}$.  Bilinearity then gives the formula.
\end{proof}

\section{What phase results for one-third anyons?}

For $m=3$ and $k=1$, winding $a$ once around $a$ multiplies the state by
$e^{2\pi i/3}$.  In an interferometer, changing the number of enclosed anyons
therefore advances the interference phase by $2\pi/3$ in the ideal abelian
model.  Real devices also include Coulomb coupling, edge reconstruction, and
decoherence; the categorical prediction concerns the topological contribution
after those effects are controlled \cite{nakamura}.

\section{What further contact should be proved?}

\begin{exercise}[A three-cocycle controls reassociation]
For $G$-graded vector spaces, multiply the associator on homogeneous degrees
$g,h,k$ by $\omega(g,h,k)\in U(1)$.  Prove that the pentagon is exactly the
$3$-cocycle equation and that changing $\omega$ by a coboundary gives a
monoidally equivalent category.
\end{exercise}

\begin{exercise}[Fusion dimensions]
If positive numbers $d_x$ satisfy
$d_xd_y=\sum_zN_{xy}^zd_z$, prove that $d$ is a positive character of the
fusion ring.  Compute it for the Fibonacci rule $\tau\otimes\tau=1\oplus\tau$.
\end{exercise}


\chapter{Chern--Simons theory turns braids into amplitudes}
\label{ch:chern-simons}

\physicalquestion{Which field theory assigns an amplitude to a braid using
only topology rather than a spacetime metric?}

\modelassumptions{Spacetime is a closed oriented three-manifold, the gauge
group is compact, and the level is chosen so the exponentiated action is gauge
invariant.  Path-integral statements are interpreted formally unless a
rigorous TQFT construction is specified.}

\section{Which action sees knots but not distances?}

For a connection $A$ on a principal $G$-bundle, the Chern--Simons action is
\[
S_{\mathrm{CS}}(A)
=\frac{k}{4\pi}\int_M
\tr\!\left(A\wedge dA+\frac23A\wedge A\wedge A\right).
\]
It uses orientation and a trace pairing but no metric.  Variation gives
$F_A=0$, so classical solutions are flat connections whose holonomy may still
be globally nontrivial.

A labeled closed worldline $K$ defines a Wilson loop
$W_R(K)=\tr_R\operatorname{Hol}_A(K)$.

\section{Why must the level be quantized?}

\begin{theorem}[Large-gauge invariance]
With the standard normalization for compact simple simply connected $G$, a
gauge transformation $g:M\to G$ changes the Chern--Simons action by
$2\pi k\,\deg(g)$.  Hence $e^{iS_{\mathrm{CS}}}$ is gauge invariant for all
gauge transformations exactly when $k\in\Z$.
\end{theorem}

\begin{proof}
Substitute $A^g=g^{-1}Ag+g^{-1}dg$ into the action.  Cyclicity of trace
separates an exact transgression term and
\[
\frac{k}{12\pi}\int_M\tr(g^{-1}dg)^3.
\]
The exact term integrates to zero on closed $M$.  Under the standard trace
normalization, the remaining integral is $2\pi k$ times the integer winding
number of $g$.  The exponential is therefore invariant for all winding
numbers precisely at integral level.
\end{proof}

\section{How do Wilson loops recover anyon data?}

In Chern--Simons TQFT, exchanging worldlines evaluates the braid morphism of a
related modular tensor category, and closing a braid evaluates its
categorical trace.  For suitable $SU(2)$ level, Wilson-loop amplitudes produce
Jones-type link invariants.  Thus fusion and braiding are not appended to the
field theory: they are the invariant data assigned to punctures and
worldlines.  The formal path integral motivates this relation; rigorous
Reshetikhin--Turaev constructions implement it categorically.
\cite{witten-jones,reshetikhin-turaev}

\section{What further contact should be proved?}

\begin{exercise}[Flat does not mean trivial]
Construct flat $U(1)$ connections with distinct holonomy on a three-torus and
compare with Chapter~\ref{ch:ab}.
\end{exercise}

\begin{exercise}[Wilson loops respect braid closure]
Assuming the ribbon-category evaluation rules, prove that isotopic closed
braids have equal amplitudes and identify where framing enters.
\end{exercise}


\chapter{Functorial theories stabilize observables}
\label{ch:generalized}

\physicalquestion{What common structure turns fusion, gluing, bundles, and
Fredholm operators into deformation-stable numerical observables?}

\modelassumptions{Categories are additive where Grothendieck completion is
used, tensor products distribute over direct sums, and generalized cohomology
is evaluated on finite CW complexes unless stated otherwise.}

\section{How are extended processes composed?}

A strong monoidal functor $F:\mathcal C\to\mathcal D$ comes with coherent
isomorphisms $F(x)\otimes F(y)\cong F(x\otimes y)$ and preserves the tensor
unit.  A three-dimensional TQFT is a symmetric monoidal functor from a bordism
category to vector spaces: disjoint union becomes tensor product and gluing
becomes composition.

For an additive monoidal category, its Grothendieck group is generated by
objects with $[x\oplus y]=[x]+[y]$ and multiplication
$[x][y]=[x\otimes y]$.  This is decategorification: fusion spaces and
morphisms are forgotten, while stable addition and multiplication remain.

\section{What survives a monoidal passage?}

\begin{theorem}[Monoidal decategorification]
An additive strong monoidal functor induces a ring homomorphism
\[
K_0(F):K_0(\mathcal C)\longrightarrow K_0(\mathcal D).
\]
If the categories are rigid and traces are preserved, categorical dimensions
and traces of closed processes are preserved as well.
\end{theorem}

\begin{proof}
Additivity gives
$[F(x\oplus y)]=[F(x)]+[F(y)]$, so the assignment $[x]\mapsto[F(x)]$
respects the defining Grothendieck relation.  The monoidal comparison gives
$[F(x\otimes y)]=[F(x)\otimes F(y)]$, hence multiplication is preserved.
Strong monoidal functors carry evaluation and coevaluation maps to duality
maps, so the composite defining categorical trace maps to the corresponding
composite in $\mathcal D$.
\end{proof}

\section{How do generalized theories unify the earlier invariants?}

A reduced generalized cohomology theory is homotopy invariant, sends cofiber
sequences to long exact sequences, and has suspension isomorphisms.  Complex
K-theory is two-periodic and admits equivalent bundle, projection, unitary,
and Fredholm models.  Those models connect the occupied Bloch bundle of
Chapter~\ref{ch:hall} to the index of Chapter~\ref{ch:index}.

The recurring pattern is now explicit:
\[
\begin{gathered}
\text{physical structure}\longrightarrow\text{functorial stable class}
\\[-2pt]
\longrightarrow\text{pairing or trace}\longrightarrow\text{observable}.
\end{gathered}
\]
Aharonov--Bohm phase, Hall conductance, Fredholm index, and anyon amplitude
use different categories, but each is protected because the middle class is
unchanged by the allowed deformation.

\section{What further contact should be proved?}

\begin{exercise}[Cohomology preserves supertraces]
Prove that cohomology is symmetric monoidal on bounded complexes of
finite-dimensional complex vector spaces and deduce the Hopf trace identity.
\end{exercise}

\begin{exercise}[Part VI synthesis]
Start with the pointed $\Z/3$ braided category of
Chapter~\ref{ch:fusion}.  Construct its braid representations, close a braid
using categorical trace, interpret the result as a Wilson-loop amplitude, and
decategorify its fusion rules.  Identify which data are lost at each passage
and verify that the $2\pi/3$ monodromy phase remains.
\end{exercise}

\begin{exercise}[Final prerequisite audit]
For one invariant from each part, list every mathematical structure on which
its definition depends.  Order those prerequisites independently of the
book's narrative and verify that every dependency appears before its first
use.  If an assumption belongs to a physical model rather than a theorem,
mark it explicitly.
\end{exercise}

\begin{quote}\textbf{Closing perspective.} Superposition demanded linear
spaces; invariant transformations demanded representations; local comparison
demanded geometry; robust global response demanded topology; matrix-level
composition demanded noncommutative analysis; and fusion with exchange
demanded categories.  At each stage, the failure of the preceding framework
selected the next coordinate-free structure.\end{quote}


\backmatter
\begin{thebibliography}{99}

\bibitem{aharonov-bohm}
Y. Aharonov and D. Bohm,
\emph{Significance of electromagnetic potentials in the quantum theory},
Physical Review 115 (1959), 485--491.

\bibitem{tonomura}
A. Tonomura et al.,
\emph{Evidence for Aharonov--Bohm effect with magnetic field completely
shielded from electron wave},
Physical Review Letters 56 (1986), 792--795.

\bibitem{choi}
M.-D. Choi,
\emph{Completely positive linear maps on complex matrices},
Linear Algebra and its Applications 10 (1975), 285--290.

\bibitem{stinespring}
W. F. Stinespring,
\emph{Positive functions on $C^*$-algebras},
Proceedings of the American Mathematical Society 6 (1955), 211--216.

\bibitem{shimanouchi}
T. Shimanouchi,
\emph{Tables of Molecular Vibrational Frequencies, Consolidated Volume I},
NSRDS--NBS 39, National Bureau of Standards, 1972.

\bibitem{tknn}
D. J. Thouless, M. Kohmoto, M. P. Nightingale, and M. den Nijs,
\emph{Quantized Hall conductance in a two-dimensional periodic potential},
Physical Review Letters 49 (1982), 405--408.

\bibitem{von-klitzing}
K. von Klitzing, G. Dorda, and M. Pepper,
\emph{New method for high-accuracy determination of the fine-structure
constant based on quantized Hall resistance},
Physical Review Letters 45 (1980), 494--497.

\bibitem{atiyah-singer}
M. F. Atiyah and I. M. Singer,
\emph{The index of elliptic operators I},
Annals of Mathematics 87 (1968), 484--530.

\bibitem{tsirelson}
B. S. Tsirelson,
\emph{Quantum generalizations of Bell's inequality},
Letters in Mathematical Physics 4 (1980), 93--100.

\bibitem{hensen}
B. Hensen et al.,
\emph{Loophole-free Bell inequality violation using electron spins separated
by 1.3 kilometres}, Nature 526 (2015), 682--686.

\bibitem{jones-index}
V. F. R. Jones,
\emph{Index for subfactors},
Inventiones Mathematicae 72 (1983), 1--25.

\bibitem{nakamura}
J. Nakamura, S. Liang, G. C. Gardner, and M. J. Manfra,
\emph{Direct observation of anyonic braiding statistics},
Nature Physics 16 (2020), 931--936.

\bibitem{witten-jones}
E. Witten,
\emph{Quantum field theory and the Jones polynomial},
Communications in Mathematical Physics 121 (1989), 351--399.

\bibitem{reshetikhin-turaev}
N. Reshetikhin and V. G. Turaev,
\emph{Invariants of 3-manifolds via link polynomials and quantum groups},
Inventiones Mathematicae 103 (1991), 547--597.

\end{thebibliography}


\end{document}
